\chapter{Trabalhos Relacionados}
\label{chap:trabalhosrelacionados}

O capítulo de Trabalhos Relacionados posiciona o estudo perante o estado da arte existente. Seu propósito não é meramente listar resumos de artigos anteriores, mas realizar uma **síntese crítica e comparativa** dos trabalhos publicados na literatura.

Ao analisar cada trabalho relacionado, o autor deve destacar:
\begin{enumerate}
    \item A abordagem ou técnica proposta pelos autores;
    \item Os resultados alcançados e suas métricas de avaliação;
    \item As principais limitações ou lacunas não resolvidas pelo estudo.
\end{enumerate}

Citações bibliográficas devem seguir rigorosamente as normas ABNT. Para citações indiretas no corpo do parágrafo, utiliza-se a forma em minúsculas como em \citeonline{akhtar2018threat}. Para citações entre parênteses ao final da afirmação, utiliza-se a forma em caixa-alta \cite{akhtar2018threat}. Em caso de referências a documentações de mercado ou recursos web, aplica-se o mesmo rigor, como em \citeonline{apple2024faceid}.

Recomenda-se finalizar este capítulo apresentando um **Quadro Comparativo** que contraste os trabalhos da literatura com a proposta deste projeto, evidenciando os diferenciais e a originalidade da contribuição.
